\bookchapter{第三十六章：手機與智能家居——使家生智}

\chapterintro{手機尚能任一日益要之事——為舉家之「遙控」。自一燈、一插座至一門鎖，智能家居方自「極客之玩」變為「家常之設」。此章為爾理清協議、生態、中樞與入門之套件，終以安全清單收之，使爾不為「智能」所反噬。}

\booksection{一、智能家居之協議：所擇至要}

購智能之器，先觀其支持何協議，此定其能為爾之機與生態所使與否。今主流之協議有四，其性各異如下：

\begin{itemize}
  \item Matter 之協議：2022 年由 CSA（連接標準聯盟）所頒，蘋果、谷歌、亞馬遜、三星諸巨擘合於其後。其志惟一——使諸品牌之智能器同其語言。2023 至 2026 年，正驟為一統之標準，支持之器可同接 Apple HomeKit、Google Home、Amazon Alexa 與米家，不復為一生態所縛。
  \item Thread 之協議：Matter 底層網絡之「管道」。耗電少、自組網（一節斷電，信號自繞他路），應之速而不滯。購帶 Thread 邊界路由器之器（如 HomePod mini、Apple TV 4K、小米中樞網關），其驗顯然更穩。
  \item Wi-Fi：帶寬大，可徑連路由器，然耗電多，器眾則易擠佔信道，宜攝像頭、掃地機之屬需傳大數據者。
  \item 藍牙（BLE）：配網簡、省電，然距短、不能自組網，宜門鎖、溫濕度計之屬小數據之器。
  \item Zigbee：老牌低耗自組網之協議，須專用之網關，器廉而穩，然跨品牌互通者劣，方為 Matter+Thread 所漸代。
\end{itemize}

\begin{tcbinfo}{一言之則}
2026 年購智能之器，先視有無「Matter over Thread」之標——此當今最經得起時驗之所擇也。舊器若僅支持 Wi-Fi 或藍牙，亦毋急棄，且用之，新購乃趨 Matter。

\end{tcbinfo}

\booksection{二、三大生態之比}

協議者，「語言」也；生態者，「圈」也。雖皆支持 Matter，日用仍須擇一生態為主，以搭自動化與聲控。三大主流生態之比如下：

\begin{center}
\begin{tabularx}{\linewidth}{|X|X|X|X|}
\hline
\textbf{生態} & \textbf{所長} & \textbf{短板} & \textbf{適用之人} \tabularnewline
\hline
Apple HomeKit & 隱私至強（以本地處理為主）、與 iPhone、Siri、快捷指令相聯甚順 & 設備偏少偏貴、須HomePod或Apple TV作中樞 & 舉家蘋果之器、重隱私與體驗之用者 \tabularnewline
米家（小米） & 器至富、性價比至高、小愛同學易用 & 生態較封閉、數據多上雲、廣告偏多 & 預算有限、欲一處配齊全屋之初學者 \tabularnewline
華為智聯（鴻蒙） & 與鴻蒙之機、平板、手錶無縫流轉、器數驟增 & 海外之器支持者少、生態尚於追趕 & 華為機之用者、重「超級終端」多器協同者 \tabularnewline
\hline
\end{tabularx}
\end{center}

\begin{tcbinfo}{擇生態之決策建議}
觀機以擇生態：iPhone 之用者先 HomeKit，華為之用者先智聯，餘安卓之用者米家最省心。恃 Matter 之協議，異生態之器亦可為一機所制，不必一次「站隊」。

\end{tcbinfo}

\booksection{三、以機為家居之中樞}

多數入門之境，機即爾「家居之中樞」——以 App 制燈光、空調、窗簾、安防，且可搭自動化之境，使家「自能運轉」：

\begin{itemize}
  \item 制器：啟 HomeKit、米家、智慧生活之 App，可遠程或本地以開關燈、調色溫、定空調之溫、啟閉窗簾、觀攝像頭之畫。
  \item 自動化之境：合「條件動作」而為之，常用三大之境： 回家之式：定位至家+天暮啟玄關之燈+啟空調+落窗簾+閉門鎖。
  \item 離家之式：末者出門（定位離家或門鎖上鎖）熄燈+關空調+閉插座+啟攝像頭+佈防。
  \item 就寢之式：每夜 23:00熄全屋主燈+落窗簾+空調切至 26 度+啟夜燈。
\end{itemize}

NFC 一碰：貼 NFC 標籤於門，機一觸即發其應之境（一觸啟門燈，再觸切就寢）。iPhone 以「快捷指令」為自動化，安卓以「錢包」或廠家之 App 皆可寫入。

\begin{tcbpractice}{試之：搭一「回家模式」}
於機智能家居 App 之「自動化」中新建一條：發之之條件擇「至某地」（家），動作依次加「啟玄關之燈」「空調設 26 度」「窗簾閉」。既存，出門一遊，歸時觀其燈自明否——初成之快，勝於新購一器。

\end{tcbpractice}

\booksection{四、入門套件推薦}

不必一次配齊，自此四件始，可蓋日用之需八成：

\begin{itemize}
  \item 智能音箱：小愛同學（米家生態，性價比高）/ HomePod mini（蘋果生態，兼為 Thread 中樞）/ Alexa（海外生態廣）。慎之：音箱定爾聲控之生態，先定生態而後購音箱。
  \item 智能燈泡：Yeelight（米家、Matter，色溫亮度可調，廉）/ Philips Hue（色佳而穩，然須 Hue 網關，稍貴）。慎之：購前核燈座之口徑（E27、B22）與需調色否。
  \item 智能插座：小米智能插座（廉、計電量）/ TP-Link Kasa（海外 Matter 支持佳）。慎之：用於大功率之器（熱水器、電暖器），必觀額定之功率，毋過載。
  \item 智能門鎖：鹿客、小米（性價比、米家相聯）/ 凱迪仕（傳統鎖具之大廠、機械之構堅實）。慎之：核其支持 Matter 或爾之生態否，留心電池之續航與應急鑰孔之所於。
\end{itemize}

\begin{tcbinfo}{選購順序建議}
先購智能音箱（定生態） 次購智能燈泡與插座（門檻低、效驗立見） 終購門鎖與攝像頭（繫乎安全，慎擇之）。一次惟增一類之器，既熟乃廣，免「購而不能用」之患。

\end{tcbinfo}

\booksection{五、安全注意事項}

家之智能器愈眾，受攻之面愈廣。下列四條，底線也，請逐條對照：

\begin{enumerate}
  \item 改其默認之密碼：攝像頭、路由器、門鎖出廠之默碼，網上盡公開矣。既得，首事即改為「長且不重複」之強碼。
  \item 及時更新固件：廠家所推之固件，多補安全之漏。於 App 中啟「自動更新」，或至少每月手查一次。
  \item 網絡之隔：置 IoT 之器於一 Wi-Fi（或訪客之網），與機、電腦之主網分之。萬一某器被破，亦難橫向脅及爾之要數據。
  \item 護其私：攝像頭啟「私域」之蔽（如寢室、廁所不攝），閉雲端二十四小時之錄影；家中攝像頭之畫面，毋長期開予第三方 App。
\end{enumerate}

\textbf{此數坑毋蹈}\par

\begin{itemize}
  \item 毋以門鎖「遠程開鎖」之功予陌生人——宜用臨時之碼，用畢即廢。
  \item 毋以攝像頭徑對牀榻、對窗之外——所洩者，爾己之私也。
  \item 毋貪廉購「白牌」無名之器——多不更新固件，為猶於家置一後門也。
\end{itemize}

\begin{tcbsummary}{本章小結}
\begin{itemize}
  \item Matter+Thread 者 2023 至 2026 年一統之向也，購新器先觀此標
  \item 三大生態各有所長：蘋果重隱私與體驗、米家重性價比、華為重多器協同
  \item 機可為中樞，搭好「回家、離家、就寢」三大之境，家遂能自運矣
  \item 入門四件：智能音箱、燈泡、插座、門鎖，依次第漸置之
  \item 改默碼、新固件、隔網絡、蔽其私——安全底線四條，不可省也
\end{itemize}

\end{tcbsummary}
